\begin{tikzpicture}
    % === Vértices principales del primer triángulo ===
    \coordinate (P)  at (0,4);  
    \coordinate (B0) at (0,0);  
    \coordinate (B1) at (3,0);  

    % === Primer triángulo ===
    \fill[LightSkyBlue3] % Color base
        (B0) -- (B1) node[midway,below] {\color{black}$K$} -- (P) -- cycle node[midway,left] {\color{black}$X$};

    % === Inicialización del patrón ===
    \coordinate (CurrOrigin) at (B1); % Origen del primer ciclo
    \coordinate (CurrTarget) at (P);  % Objetivo del primer ciclo

    % === Bucle de iteración geométrica ===
    % Generación de la secuencia con altura decreciente
    \foreach \i [evaluate=\i as \h using 1.5/(2^(\i-2))] in {2,3,4,...,8} {
        % Cálculo de la base actual
        \coordinate (CurrBase) at ($(CurrOrigin)!4cm!(CurrTarget)$);
        
        % Cálculo del punto perpendicular hacia afuera
        \coordinate (CurrP) at ($(CurrOrigin)!\h cm!-90:(CurrBase)$);

        % Asignación de color alternante
        \ifodd\i
            \def\tcolor{LightSkyBlue3} % Color para iteraciones impares
        \else
            \def\tcolor{DodgerBlue4}   % Color para iteraciones pares
        \fi

        % Trazado del triángulo
        \fill[fill=\tcolor]
            (CurrOrigin) -- (CurrBase) -- (CurrP) -- cycle;

        % Actualización de puntos para el siguiente ciclo
        \coordinate (CurrTarget) at (CurrBase);
        \coordinate (CurrOrigin) at (CurrP);
    }
\end{tikzpicture}